\begin{Huge}
    Medieninformatik-FAQ
\end{Huge}

\begin{block}{H\"aufig gestellte Fragen zum Studium}

\question{Lernt man im Studium, wie man programmiert?}{
    Ja, aber auf einer sehr eigenst\"andigen Basis. Man bekommt einen \"Uberblick \"uber die Sprache(n) und \"uber dahinterstehende Programmierkonzepte. Praktische Programmiererfahrung muss man sich aber gr\"o"stenteils selbst erarbeiten.
}

\question{Welche Programmiersprachen macht man da so?}{
    Das ist von Prof zu Prof und von Veranstaltung zu Veranstaltung unterschiedlich. In den ersten beiden Semestern lernt man meistens Racket und Java.
}

\question{Muss man programmieren k\"onnen, um das Studium anzufangen?}{
    Nein. Die Vorlesung beginnt absolut bei 0, um allen den Einstieg zu erm\"oglichen.
}

\question{Muss man gut in Mathe sein?}{
    Man muss kein Mathe-Genie sein, man sollte Mathe aber nicht hassen. Es ist gerade am Anfang viel Mathe. Auch sollte man sich nicht davon t\"auschen lassen, dass es in der Medieninformatik eine Mathe-Vorlesung weniger gibt wie in der reinen Informatik. Denn auch andere Vorlesungen wie Graphische Datenverarbeitung k\"onnen sehr Mathe-lastig sein.
}

\question{Ich besch\"aftige mich gerne mit Design und so, kann ich Medieninformatik studieren?}{
    Medieninformatik ist in T\"ubingen vor allem eines: Informatik. Gestaltung von Nutzeroberfl\"achen geh\"ort zwar zum Studium dazu, im regul\"aren Studium gibt es aber quasi keine Module, in denen man praktisch designt. Viel mehr geht es um theoretische Grundlagen, insbesondere von der Informatik-Seite her.
}

\question{Wie ist die Frauenquote so?}{
    35,2\%.
}

\question{Lerne ich im Studium, wie man Webseiten baut?}{
    In einer Veranstaltung, ja. Die ersten ca. 3 Semester haben viel mit Web und User-Interfaces zu tun, danach verschiebt sich der Schwerpunkt je nach pers\"onlicher Pr\"aferenz.
}

\question{Lerne ich, wie man Computerspiele baut?}{
    Die M\"oglichkeit besteht, ist aber nicht grundlegender Teil des Studiums und kommt, wenn \"uberhaupt, erst im dritten Studienjahr oder sp\"ater.
}

\question{Was arbeitet man danach so?}{
    Alle Bereiche der IT-Branche, insbesondere Webentwicklung, Entwicklung von Computerspielen, in der Filmindustrie, Automobilbranche und Medizintechnik.
}

\question{Gibt es Praktika?}{
    Im normalen Studienverlauf ist kein berufsorientiertes Praktikum vorgesehen, viele arbeiten aber parallel als Werksstudent*in oder man macht ein Kurzpraktikum in den Semesterferien.
}

\question{Kann man ein Auslandssemester machen?}{
    Klar, geht immer. T\"ubingen nimmt am ERASMUS-Programm teil, die Organisation ist aber langwierig und man sollte sich fr\"uh drum k\"ummern.
}

\question{Wie ist da so der NC?}{
    Die letzten 5 Jahre gab es keinen NC. Das muss aber theoretisch nicht immer so sein.
}

\end{block}
